% ============================================================
% sec/04track_analysis.tex
% Keep ONLY the subsection matching the approved track.
% Delete the other two before submission.
% ============================================================
\section{Technical Analysis of the Selected Track}
\label{sec:trackanalysis}

\todo{One sentence naming the approved track and the date of approval.}

% ------------------------------------------------------------
% TRACK A --- Classical machine learning (tabular)
% ------------------------------------------------------------
\subsection{Exploratory Data Analysis}
\label{subsec:eda}
\todo{Number of records, number of features, target distribution,
temporal coverage.}

\subsubsection{Data Dictionary}
\begin{table}[!t]
\centering
\caption{Data dictionary.}
\label{tab:datadict}
\footnotesize
\begin{tabularx}{\columnwidth}{@{}P{1.8cm} P{1.4cm} Y@{}}
\toprule
\textbf{Variable} & \textbf{Type} & \textbf{Meaning / range / missing \%}\\
\midrule
\todo{var} & \todo{type} & \todo{description}\\
\bottomrule
\end{tabularx}
\end{table}

\subsubsection{Data Quality and Imputation}
\todo{Missing value pattern and the justified imputation strategy.}

\subsubsection{Outlier Detection}
\todo{Method, threshold, and what is done with the detected outliers.}

\subsubsection{Class Imbalance}
\todo{Ratio, and the strategy chosen with its justification.}

% ------------------------------------------------------------
% TRACK B --- Deep learning (NLP / vision)
% ------------------------------------------------------------
\subsection{Corpus / Image Set Analysis}
\label{subsec:corpus}
\todo{Size, class distribution, sequence length or resolution
distribution, documented biases of the source.}

\subsubsection{Preprocessing and Tokenization}
\todo{Pipeline and justification.}

\subsubsection{Data Augmentation}
\todo{Transformations, why each one is valid for this domain, and which
were discarded as label-destroying.}

\subsubsection{Hardware and Memory Considerations}
\todo{Model size, batch size, precision, expected GPU memory, and
estimated training time on Google Colab Pro.}

% ------------------------------------------------------------
% TRACK C --- Reinforcement learning
% ------------------------------------------------------------
\subsection{MDP Formalization}
\label{subsec:mdp}
The task is modeled as a Markov decision process
$\mathcal{M} = \langle \mathcal{S}, \mathcal{A}, P, R, \gamma \rangle$.

\subsubsection{\texorpdfstring{State Space $\mathcal{S}$}{State Space S}}
\todo{Exact variables, ranges, normalization, and whether the
environment is fully or partially observable.}

\subsubsection{\texorpdfstring{Action Space $\mathcal{A}$}{Action Space A}}
\todo{Discrete or continuous, dimensionality, bounds.}

\subsubsection{\texorpdfstring{Transition Dynamics $P(s'\mid s,a)$}{Transition Dynamics}}
\todo{Simulator, its version, stochasticity sources, timestep duration,
and episode termination conditions.}

\subsubsection{\texorpdfstring{Reward Function $R(s,a,s')$}{Reward Function}}
\begin{equation}
\label{eq:reward}
R(s,a,s') = \todo{reward expression}
\end{equation}
\todo{Justify every term and every weight. State which reward designs
were discarded and which degenerate behavior each one would induce.}

\subsubsection{\texorpdfstring{Discount Factor $\gamma$}{Discount Factor}}
\todo{Value and justification with respect to the effective horizon.}

\subsubsection{Environment Verification}
\todo{Smoke test executed, random-policy return, and confirmation that
the environment runs on Google Colab Pro.}